% W4i46_Compiled.tex
% DFD 17-Agent Research Programme --- Iteration 46 (i46)
% ALL 17 AGENTS: Master Compiled Document
% Theme: Phase 0A Execution Plan, Birefringence CS Assessment, JUNO First Results, Euclid 88 Days
% Date: 2026-03-23

\documentclass[11pt,a4paper]{article}
\usepackage[utf8]{inputenc}
\usepackage{amsmath,amssymb,amsfonts,amsthm}
\usepackage{hyperref}
\usepackage{booktabs}
\usepackage{longtable}
\usepackage{geometry}
\usepackage{xcolor}
\usepackage{tcolorbox}
\geometry{margin=2.5cm}

\newtheorem{theorem}{Theorem}
\newtheorem{proposition}{Proposition}
\newtheorem{lemma}{Lemma}
\newtheorem{corollary}{Corollary}
\newtheorem{definition}{Definition}
\newtheorem{remark}{Remark}

\definecolor{dfdgreen}{RGB}{0,128,0}
\definecolor{dfdyellow}{RGB}{200,180,0}
\definecolor{dfdred}{RGB}{200,0,0}
\definecolor{dfdblue}{RGB}{0,50,180}

\newcommand{\GREEN}{\textcolor{dfdgreen}{\textbf{GREEN}}}
\newcommand{\YELLOW}{\textcolor{dfdyellow}{\textbf{YELLOW}}}
\newcommand{\RED}{\textcolor{dfdred}{\textbf{RED}}}
\newcommand{\NEW}{\textcolor{dfdblue}{\textbf{NEW}}}
\newcommand{\RESOLVED}{\textcolor{dfdgreen}{\textbf{RESOLVED}}}
\newcommand{\Tr}{\operatorname{Tr}}
\newcommand{\CP}{\mathbb{C}P}

\title{DFD i46: Compiled Master Document --- All 17 Agents\\[6pt]
\large Iteration 15/20 of Quantum Gravity Cycle (i32--i51)\\[6pt]
\normalsize Phase 0A Plan $\cdot$ Birefringence CS Analysis $\cdot$ JUNO First Results $\cdot$ Scorecard 61/5/4}
\author{DFD 17-Agent Research Programme}
\date{2026-03-23}

\begin{document}
\maketitle

\begin{abstract}
Iteration 46 advances the programme on three fronts. First, Agent 14 has produced a fully specified 12-hour, 4-step execution plan for Phase 0A (background-only SymBoltz computation with parameterized distance-bias), accompanied by a theoretical framework for Phase 0B ($\psi$ perturbation theory derivation). Second, Agent 11 has performed the birefringence Chern--Simons risk assessment: DFD's scalar $\psi$-field cannot produce birefringence directly, but three accommodation options exist, with the most elegant being evolution of the Chern--Simons sector that already determines $\alpha$. Third, JUNO has published its first physics results with world-leading precision on $\Delta m^2_{21}$ and $\theta_{12}$, and the combination with T2K+NOvA shows a $2.2\sigma$ preference for normal ordering, consistent with DFD. Additionally: the b-quark exponent is resolved ($n_b = 0$ adopted), the PRL impediment is partially resolved (does not block DFD internal torsion), ET is now a three-site race (decision 2027), Gaia DR4 is confirmed for December 2026, and the Euclid pre-registration deadline is 88 days away. Scorecard remains at 61 GREEN, 5 YELLOW, 4 RED --- unchanged for the fourth consecutive iteration.
\end{abstract}

\tableofcontents
\newpage

%=============================================================================
\part{Executive Summary}
%=============================================================================

\section{i46 vs.\ i45 Comparison}

\begin{center}
\begin{tabular}{lll}
\toprule
\textbf{Aspect} & \textbf{i45} & \textbf{i46} \\
\midrule
Critical resolutions & 1 (Phase 0 blocker) & 2 (b-quark, PRL impediment partial) \\
Theory upgrades & 0 & 1 ($\psi$ perturbation framework) \\
New literature & $85+$ papers & $85+$ papers ($\sim$30 new) \\
Scorecard & 61/5/4 & 61/5/4 \\
Phase 0 & 0A/0B split proposed & \textbf{0A fully specified; 0B framework outlined} \\
Key external & GWTC-4, Moriond 2026 & JUNO first results, ET three-site, Gaia DR4 date \\
Grade changes & 2 up, 2 down & 0 changes \\
Birefringence & Escalated to MEDIUM-HIGH & \textbf{CS accommodation analysis performed} \\
\bottomrule
\end{tabular}
\end{center}

\section{Top 5 New Developments}

\begin{enumerate}
\item \textbf{Phase 0A execution plan fully specified}: 12 hours, 4 steps, explicit parameter choices. Phase 0B theoretical framework for $\psi$ perturbation theory outlined. Key literature identified (K-essence Boltzmann dynamics, $F(T)$ gravity in Boltzmann codes, SymBoltz custom submodels).

\item \textbf{Birefringence Chern--Simons assessment}: Three options analyzed. Option 1 (independent source, trivial but inelegant), Option 2 ($\psi F \tilde{F}$ coupling, requires pseudoscalar extension), Option 3 (CS sector evolution, elegant but requires non-integer CS levels). Anisotropic birefringence measured consistent with zero (arXiv:2504.13154), which is mildly encouraging for DFD.

\item \textbf{JUNO first results published}: 59.1 days of data achieve world-leading $\Delta m^2_{21}$ and $\theta_{12}$ precision. JUNO + T2K + NOvA: normal ordering preferred at $2.2\sigma$. DFD predicts normal ordering. The $3\sigma$ milestone expected by 2027--2028.

\item \textbf{b-quark resolved and PRL impediment partially resolved}: $n_b = 0$ adopted as baseline (0.3\% PDG match). PRL impediment applies to spacetime torsion, not DFD's internal fibre torsion --- does not kill the spectral torsion route.

\item \textbf{Euclid/Gaia timelines sharpened}: Euclid DR1 confirmed 21 October 2026. Gaia DR4 no earlier than December 2026. Euclid pre-registration deadline 20 June 2026 (88 days). Higher-order WL statistics outperform two-point by $2.5\times$ (Euclid Prep.\ LXXXV).
\end{enumerate}


%=============================================================================
\part{Agent Grades: All 17}
%=============================================================================

\begin{center}
\begin{longtable}{clp{9cm}}
\toprule
\textbf{Agent} & \textbf{Grade} & \textbf{Key Result} \\
\midrule
\endfirsthead
\toprule
\textbf{Agent} & \textbf{Grade} & \textbf{Key Result} \\
\midrule
\endhead
1 ($\alpha$) & B+ & $K_\alpha$ confirmed published. Quasar convergence supports DFD. \\
2 (Topology) & B+ & PRL impediment partially resolved: does not kill internal torsion route. \\
3 (Masses) & A- & b-quark $n_b = 0$ adopted. Theorem K.4 write-up specified. \\
4 (Spin$^c$) & B & $k_f$ gap remains. Conceptual link to Phase 0B noted. \\
5 (Anomaly) & C+ & Mixed anomaly deferred 6 iterations. JUNO NO at $2.2\sigma$. \\
6 (Dead patents) & C+ & Patent landscape unchanged. \\
7 (ET/CE) & A- & ET three-site race. IR1 Sept--Oct 2026. GW background limits. \\
8 (Laboratory) & A- & Th-229 Nature published. DM detection method proposed. \\
9 (Particle) & B+ & Moriond 2026 fully digested. $H \to Z\gamma$ at $2.5\sigma$. \\
10 (Astro) & B+ & Gaia DR4 confirmed December 2026. \\
11 (CMB) & B+ & \textbf{Birefringence CS assessment: three options analyzed.} \\
12 (LSS) & B+ & Euclid DR1 HO-WL $2.5\times$ improvement. 88-day countdown. \\
13 (Dark) & B & JUNO first results: NO at $2.2\sigma$. $\Sigma m_\nu$ Feldman--Cousins tension. \\
14 (Numerics) & B & Phase 0A plan fully specified. Phase 0B framework outlined. Unexecuted. \\
15 (Scorecard) & B & 61/5/4 frozen for 4 iterations. \\
16 (Cross-check) & A- & All claims verified. Phase 0 and mixed anomaly process failures flagged. \\
17 (Synthesis) & A- & Honest trajectory assessment. Five iterations remaining. \\
\bottomrule
\end{longtable}
\end{center}

\textbf{Mean grade}: B+ (unchanged from i45).

\textbf{Grade changes} (i45 $\to$ i46): \textbf{None.} All agents hold their i45 grades. The absence of grade movement reflects the programme's holding pattern: significant intellectual work (CS analysis, Phase 0B framework, b-quark resolution) but no numerical results or scorecard movement.


%=============================================================================
\part{Scorecard: 70 Predictions}
%=============================================================================

\section{Summary}

\begin{center}
\begin{tabular}{lccc}
\toprule
& \textbf{\GREEN} & \textbf{\YELLOW} & \textbf{\RED} \\
\midrule
i45 & 61 & 5 & 4 \\
i46 & 61 & 5 & 4 \\
\textbf{Change} & \textbf{0} & \textbf{0} & \textbf{0} \\
\bottomrule
\end{tabular}
\end{center}

Fourth consecutive iteration at 61/5/4. Promotion blocked by Phase 0 (for $A_L$), Gaia DR4 (for EFE), and JUNO (for $\Sigma m_\nu$).

\section{RED Predictions (4)}

\begin{enumerate}
\item GW echo amplitude ($10^{-41}$): untestable
\item BMV enhancement ($2200\times$ QED): awaiting data
\item Lepton universality: anomalies dissolved, pending Run 3
\item $|\lambda-1|$ direct: no experiment sensitive
\end{enumerate}

\section{YELLOW Predictions (5)}

\begin{enumerate}
\item $\dot{\alpha}/\alpha$ at $10^{-19}$/yr: Th-229 by $\sim$2030 (on track)
\item EFE in galaxies: converging to GR, Gaia DR4 December 2026
\item $A_L$ lensing: $1.025 \pm 0.017$ vs.\ DFD $1.00$--$1.05$ (Phase 0A needed)
\item $\Sigma m_\nu$: $< 0.064$ eV vs.\ DFD $0.061$ eV (Feldman--Cousins tension emerging)
\item Dynamical dark energy: DESI $2.8$--$4.2\sigma$ (debated)
\end{enumerate}


%=============================================================================
\part{Phase 0: Critical Path Update}
%=============================================================================

\section{Phase 0A: Fully Specified, Awaiting Execution}

\begin{tcolorbox}[colback=green!5!white, colframe=green!60!black,
  title={Phase 0A: 12-Hour Execution Plan --- Ready}]
\begin{enumerate}
\item \textbf{Hours 0--2}: Install SymBoltz.jl v1.0.0. Verify $\Lambda$CDM baseline.
\item \textbf{Hours 2--5}: Implement distance-bias $D_L^{\text{DFD}} = D_L^{\Lambda\text{CDM}} \cdot e^{\Delta\psi(z)}$ with $\Delta\psi(z) = a_1 z + a_2 z^2 + a_3 z^3$, coefficients from DFD field equations at background level.
\item \textbf{Hours 5--8}: Compute $A_L^{\text{DFD}}$ (background-only) and BAO distance ratios.
\item \textbf{Hours 8--12}: Generate Euclid pre-registration numbers. Fisher forecast.
\end{enumerate}

\textbf{Status}: Plan fully specified. Not executed.
\end{tcolorbox}

\section{Phase 0B: Framework Outlined}

\begin{tcolorbox}[colback=yellow!5!white, colframe=yellow!60!black,
  title={Phase 0B: $\psi$ Perturbation Theory --- Framework}]
\textbf{Five parameters to determine}:
\begin{enumerate}
\item $\psi$ effective mass $m_\psi$ (from potential second derivative)
\item Sound speed $c_s^2$ ($= 1$ for canonical scalar)
\item Anisotropic stress $\pi_\psi$ ($= 0$ for canonical scalar)
\item Coupling to photon/baryon perturbations (via refractive index)
\item Initial conditions for $\delta\psi$ at recombination
\end{enumerate}

\textbf{Key literature for implementation}:
\begin{itemize}
\item arXiv:2602.19576 --- K-essence Boltzmann dynamics (modified mass-shell, geodesics)
\item arXiv:2512.16404 --- $F(T)$ gravity in Boltzmann codes (framework)
\item SymBoltz.jl --- Custom submodel support (published A\&A March 2026)
\end{itemize}

\textbf{Timeline}: 1--2 months from start of focused work.

\textbf{Status}: Framework outlined. Derivation not started.
\end{tcolorbox}


%=============================================================================
\part{Birefringence: Risk Assessment and Accommodation Options}
%=============================================================================

\section{Current Data}

\begin{itemize}
\item Planck PR4 isotropic: $\beta = (0.30 \pm 0.11)^\circ$ ($\sim 3\sigma$).
\item PRL January 2026: angle may be larger with new uncertainty treatment.
\item Anisotropic: $0.42^{+0.40}_{-0.34} \times 10^{-4}$, consistent with zero.
\item SO + LiteBIRD will be decisive ($\sim$2027--2028).
\end{itemize}

\section{DFD Accommodation Options}

\begin{center}
\begin{tabular}{lp{5cm}cc}
\toprule
\textbf{Option} & \textbf{Description} & \textbf{Compatibility} & \textbf{Economy} \\
\midrule
1. Independent source & Birefringence from ALPs/B-fields, unrelated to DFD & Full & Poor \\
2. $\psi F \tilde{F}$ coupling & Scalar $\psi$ couples to photon dual & Requires extension & Medium \\
3. CS sector evolution & Chern--Simons sector that gives $\alpha$ also rotates polarization & Requires non-integer CS & Best \\
\bottomrule
\end{tabular}
\end{center}

\textbf{Assessment}: Risk level MEDIUM-HIGH. Not yet at crisis level ($< 5\sigma$). Options exist. Monitoring continues.


%=============================================================================
\part{Falsification Risk Dashboard}
%=============================================================================

\begin{center}
\begin{tabular}{lcccl}
\toprule
\textbf{Observable} & \textbf{DFD} & \textbf{Data} & \textbf{Risk} & \textbf{i46 Change} \\
\midrule
$\Sigma m_\nu$ & 0.061 eV & $< 0.064$ eV & \YELLOW & FC tension $3\sigma$ (noted) \\
$A_L$ & 1.00--1.05 & $1.025 \pm 0.017$ & \GREEN & No change \\
$R_{31}$ & TBD & --- & \YELLOW & No change \\
Birefringence & None & $\geq 0.3^\circ$ & \RED & \textbf{CS analysis performed} \\
GW lensing & Never lensed & None & \GREEN & No change \\
$D_L^{\text{EM}}/D_L^{\text{GW}}$ & $\sim 1.31$ & N/A & \GREEN & No change \\
Mass ordering & Normal & $2.2\sigma$ NO & \GREEN & \textbf{JUNO trending right} \\
\bottomrule
\end{tabular}
\end{center}

\textbf{Birefringence remains the highest-risk item}. $\Sigma m_\nu$ Feldman--Cousins tension is an emerging secondary concern.


%=============================================================================
\part{New Literature: Complete List ($\sim$85 papers across 17 agents)}
%=============================================================================

\section{Theory (Agents 1--5)}

\begin{enumerate}
\item Nature Communications (2025) --- $K_\alpha$ sensitivity for Th-229. \NEW.
\item MNRAS 528, 6550 (2024) --- Quasar $\alpha$ measurement convergence. \NEW.
\item JHEP (February 2025) --- Spectral torsion for Connes type operator. \NEW.
\item Uppsala thesis (2025) --- Spin, Spectral Geometry and Gravity. \NEW.
\item JHEP (February 2026) --- LEFT two-loop RGEs. \NEW.
\item arXiv:2511.14593 (November 2025) --- JUNO first reactor oscillation measurement. \NEW.
\item arXiv:2601.09791 (January 2026) --- Lessons from first JUNO results. \NEW.
\item arXiv:2603.17181 (March 2026) --- JUNO + long-baseline synergy. \NEW.
\item Phys.\ Rev.\ D 111 (2025) --- Neutrino mass cosmological limits, beyond-$\Lambda$CDM. \NEW.
\item arXiv:2508.16238 --- Neutrino masses in second-order CPL DE. \NEW.
\item arXiv:2506.18328 (June 2025) --- Comprehensive $\alpha$ review. (continued)
\item arXiv:2601.22273 (January 2026) --- Lattice QCD excited-state uncertainties. (continued)
\item arXiv:2503.14744 --- DESI DR2 + Planck $\Sigma m_\nu$. (continued)
\item T2K + NOvA (Nature, October 2025) --- Joint analysis. (continued)
\end{enumerate}

\section{Observables (Agents 8--14)}

\begin{enumerate}
\item Nature (January 2026) --- Th-229 frequency reproducibility (published). \NEW.
\item Dark matter detection via Th-229 nuclear clock (July 2025). \NEW.
\item ATLAS Moriond 2026 full summary --- including $H \to Z\gamma$ at $2.5\sigma$. \NEW.
\item ESA Gaia DR4 scenario --- December 2026. \NEW.
\item arXiv:2510.17746 --- ML identification of Gaia astrometric exoplanet orbits. \NEW.
\item arXiv:2504.13154 (June 2025) --- Anisotropic birefringence from B-mode, null. \NEW.
\item arXiv:2602.12019 (February 2026) --- Simulation-based birefringence inference. \NEW.
\item arXiv:2512.19102 (December 2025) --- SO SAT polarization calibration. \NEW.
\item Euclid Prep.\ LXXXV (A\&A, March 2026) --- DR1 higher-order WL. \NEW{} (published).
\item arXiv:2507.07629 --- Euclid ERO WL pipeline validation. \NEW.
\item arXiv:2508.16238 --- Second-order CPL neutrino mass bounds. \NEW.
\item arXiv:2602.19576 (February 2026) --- K-essence Boltzmann dynamics. \NEW.
\item arXiv:2512.16404 (December 2025) --- $F(T)$ gravity in Boltzmann codes. \NEW.
\item Qu et al.\ (PRL 136, 2026) --- $A_L = 1.025 \pm 0.017$. (continued)
\item SymBoltz.jl (A\&A, March 2026) --- now published. (continued)
\end{enumerate}

\section{Cosmology and GW (Agents 6, 7, 15, 16, 17)}

\begin{enumerate}
\item ET EMR student team (March 2026) --- Bid book preparation. \NEW.
\item ET-SUnLab tender published (February 2026) --- EU Official Journal. \NEW.
\item ET INGV Mamone station (March 2026) --- Geophysical characterization. \NEW.
\item Sardinia--Saxony cooperation (January 2026) --- Declaration of intent. \NEW.
\item arXiv:2508.20721 --- O4a isotropic GW background limits. \NEW.
\item LVK IR1 planning --- September--October 2026. \NEW.
\item GWTC-4.0 (arXiv:2508.18080) --- 218 events. (continued)
\item ET-SUnLab tender (February 2026) --- (continued)
\end{enumerate}


%=============================================================================
\part{Action Items for i47}
%=============================================================================

\begin{center}
\begin{tabular}{clcc}
\toprule
\textbf{Priority} & \textbf{Action} & \textbf{Owner} & \textbf{Blocks} \\
\midrule
\textbf{CRITICAL} & \textbf{Execute Phase 0A} (not plan --- execute) & Agent 14 & Euclid, $A_L$, scorecard \\
\textbf{HIGH} & Submit qualitative Euclid pre-registration & Agent 12 & Phase 0A \\
\textbf{HIGH} & Begin Phase 0B derivation (5 parameters) & Agent 14 & Full predictions \\
HIGH & Edit mass paper file with Theorem K.4 & Agent 3 & Publication \\
MEDIUM & Decide: attempt mixed anomaly or drop & Agent 5 & Agenda clarity \\
MEDIUM & Track birefringence + SO calibration & Agent 11 & Falsification risk \\
MEDIUM & Monitor JUNO mass ordering accumulation & Agent 13 & $\Sigma m_\nu$ \\
MEDIUM & Track $\Sigma m_\nu$ Feldman--Cousins development & Agent 5/13 & Prediction viability \\
LOW & Quantitative PRL impediment on $k_f$ & Agent 2 & Mass exponents \\
\bottomrule
\end{tabular}
\end{center}

\textbf{Deadline reminder}: Euclid pre-registration by 20 June 2026 (\textbf{88 days}). LHC Run 3 ends June 2026. ET-SUnLab tender deadline 16 April 2026 (24 days). LVK IR1 September--October 2026. Gaia DR4 December 2026.

\textbf{Programme status}: Iteration 15/20 (75\% complete). Five iterations remain. Phase 0A execution is the single most important action for the remaining programme.

\end{document}
